
\section{Definir el indicador}
Un indicador necesita población, evento, exposición, periodo y reglas. Cinco retiradas no equivalen necesariamente a cinco fallos confirmados. Una tasa dividida por horas de otra población carece de interpretación. Se identifican unidades repetidas, modificaciones y utilización.

Bajo Poisson homogéneo, $\widehat\lambda=n/T$ estima tasa; con supervivencia exponencial, $R(t)=\exp(-\lambda t)$. Proceso de eventos reparables y distribución de vida de una pieza se distinguen. MTBF no es límite de vida ni periodo libre de fallos garantizado. P32 explica por qué cero eventos no prueba tasa cero.
\section{Weibull y observación incompleta}
Con $R(t)=\exp[-(t/\eta)^\beta]$, la tasa instantánea es $h(t)=(\beta/\eta)(t/\eta)^{\beta-1}$. En el modelo crece con edad si $\beta>1$, es constante con $\beta=1$ y decrece con $\beta<1$. La forma estadística no identifica automáticamente mecanismo físico ni una curva de bañera universal.

Las unidades aún operativas aportan censura, no fallos en la fecha de cierre. Ignorarlas puede sesgar la estimación. Censura y pérdida de historial también difieren. Supervivencia condicional usa $R(t_0+\Delta t)/R(t_0)$, generalmente distinta de $R(\Delta t)$ salvo el caso exponencial.
\section{Recursos, repuestos y riesgo}
Planificar exige coincidencia de instalación, personal, útil, datos y repuesto dentro de vencimientos. Agrupar tareas puede reducir paradas y adelantar consumo de intervalos o elevar complejidad. La planificación no desplaza una obligación para mejorar disponibilidad.

Un repuesto físicamente presente puede carecer de elegibilidad o preservación. Suministro, reparación e incertidumbre condicionan stock y parada. Coste de ciclo de vida incluye más que adquisición. Las matrices cualitativas ayudan a priorizar, sin convertir escalas ordinales en probabilidades precisas. Se conserva evidencia y se comprueba eficacia de medidas dentro del sistema pertinente.
\repaso{¿Reducir MTTR excluyendo esperas demuestra mejora operacional?}{Puede haber cambiado solo la definición. Se comparan tiempos de igual alcance y se comprueba si las paradas reales disminuyeron.}
